\documentclass[10pt,a4paper]{amsart}
\usepackage[margin=19mm]{geometry}
\usepackage{amsmath,amssymb,mathtools}
\usepackage[scr=rsfs,cal=euler]{mathalfa}
\usepackage{xfrac}
\usepackage[unicode,hidelinks,bookmarks=false]{hyperref}
\newcommand{\ie}{i.e.\ }
\newcommand{\cf}{cf.\ }
\newcommand{\resp}{respectively\ }
\newcommand{\Subsection}{\subsection}
\providecommand{\nobreakdash}{\nobreak\hbox{-}}
\setlength{\emergencystretch}{2em}
\multlinegap0pt
\usepackage{float}
\usepackage{dsfont}
\usepackage{amssymb}
\newtheorem{thm}{Theorem}
\newcommand{\myseries}[3]{#1\leq #2\leq #3}
\newcommand{\myvec}[3]{#1=(#2_1,#2_2,\ldots,#2_#3)}
\newcommand{\norm}[2]{\lVert #1 \rVert_#2}
\newcommand{\prodifs}[1]{#1_{i_1i_2\cdots i_m}}
\newcommand{\prodpoint}[2]{#1_{1,#2_1},#1_{2,#2_2},\ldots,#1_{m,#2_m}}
\newcommand{\profunc}[1]{#1_1\times#1_2\times\cdots\times#1_m}
\title{Study of attractors and fractal functions on the product spaces and Dimensional aspects}
\author{Alamgir Hossain}
\date{Source: arXiv:2603.12890v1 (CC BY 4.0)}
\begin{document}
\maketitle

\noindent\emph{Source and licence.} Study of attractors and fractal functions on the product spaces and Dimensional aspects. Version arXiv:2603.12890v1 (CC BY 4.0), distributed by arXiv under the Creative Commons Attribution 4.0 International licence, http://creativecommons.org/licenses/by/4.0/, as recorded in the arXiv licence metadata for this exact version and shown on the abstract page https://arxiv.org/abs/2603.12890v1. Full licence notice: \texttt{licenses/arxiv-2603.12890-CC-BY-4.0.txt}.\par\smallskip

\noindent\textit{Editorial scope.} Counted unit: the article's thm graphofcpf (source0.tex lines 656--658). The article's complete original proof of it is delivered with it (source0.tex lines 659--737, one \texttt{proof} environment of 79 lines). No mathematics is written, repaired or re-derived: the statement and the proof are the article's own lines, in the article's own notation, and no retained line is substituted or re-flowed.  Retained uncounted as the source-local context the proof needs: lines 536--538; lines 549--552; lines 572--590.  Nothing was omitted from any retained span, so the package carries no omission record.  No retained line contains a citation, so the wrapper carries no bibliography and no bibliography entry.  No retained line contains a figure environment, an image-inclusion command or drawing code, so no image is delivered and the package declares no asset dependency.  Editorial formatting only, in addition: an AMS \texttt{amsart} preamble that restates the article's own declarations and macros used by the retained lines, namely the environments \texttt{thm} on separate counters, together with the standard packages those lines need.  The retained environments print in this document as Theorem 1 in order of appearance, whereas the article prints the same items with the numbering of its own full document.  The complete native source of the article as distributed in the arXiv e-print for this exact version is bundled unchanged as \texttt{sources/196300/source0.tex}.
	\begin{equation}
		(T_kf)(x)=F_{k,i_k}(L_{k,i_k}^{-1}(x),f(L_{k,i_k}^{-1}(x)))\quad\mbox{for}~x\in I_{k,i_k};~\myseries{1}{i_k}{N_k}. \label{prodrbop}
	\end{equation}  

	\begin{equation}
		L_{k,i_k}(x_{k,j_k})=x_{k,\sigma(i_k,j_k)},\quad F_{k,i_k}(x_{k,j_k},y_{k,j_k})=y_{k,\sigma(i_k,j_k)}
		\label{contpoint}
	\end{equation}

	Further, the maps $\prodifs{F}:I\times\mathbb{R}^m\rightarrow\mathbb{R}^m$ are defined by 
	\begin{equation}
		\prodifs{F}(\vec{x},\vec{y})=(F_{1,i_1}(x_1,y_1),F_{2,i_2}(x_2,y_2),\ldots,F_{m,i_m}(x_m,y_m)) \label{prod of F}
	\end{equation}
	for $\myvec{\vec{x}}{x}{m}\in I,~\myvec{\vec{y}}{y}{m}\in\mathbb{R}^m.$ Hence from (\ref{contpoint}), we get 
	\begin{equation}
		\prodifs{F}((\prodpoint{x}{j}),(\prodpoint{y}{j}))=(y_{1,\sigma(i_1,j_1)},y_{2,\sigma(i_2,j_2)},\ldots,y_{m,\sigma(i_m,j_m)})
		\label{prodofF}
	\end{equation}
	for all $i_k\in \mathbb{N}_{N_k},~j_k\in \partial\mathbb{N}_{N_k,0},~ k\in\mathbb{N}_m$ and 
	\begin{equation}
		\left. \begin{array}{ll}
			\lVert \prodifs{F}(\vec{x},\vec{y})-\prodifs{F}(\vec{x'},\vec{y})\rVert_2\leq |\prodifs{a}|_{\infty}\lVert \vec{x}-\vec{x'}\rVert_2,\quad\mbox{for all}~ x.x'\in I\\
			\lVert \prodifs{F}(\vec{x},\vec{y})-\prodifs{F}(\vec{x},\vec{y'})\rVert_2\leq |\prodifs{b}|_{\infty}\lVert \vec{y}-\vec{y'}\rVert_2,\quad\mbox{for all}~ y,y'\in\mathbb{R}^m.
		\end{array} \right. 
		\label{cont.of F}
	\end{equation}
	where $|\prodifs{a}|_{\infty}=\max\{|a_{k,i_k}|:\myseries{1}{i_k}{N_k};~\myseries{1}{k}{m}\}<1$ and  $|\prodifs{b}|_{\infty}=\max\{|b_{k,i_k}|:\myseries{1}{i_k}{N_k};~\myseries{1}{k}{m}\}<1$.\\
	Finally, the maps $\prodifs{W}:I\times\mathbb{R}^m\rightarrow\prodifs{I}\times\mathbb{R}^m$ are defined by 

	\begin{thm}
		The attractor $G$ of the IFS $\{(I\times\mathbb{R}^m;\prodifs{W}):~\myseries{1}{i_k}{N_k};\myseries{1}{k}{m}\}$, associated with the data set $\{((\prodpoint{x}{i}),(\prodpoint{y}{i}))\in I\times\mathbb{R}^m:~\myseries{1}{i_k}{N_k};\myseries{1}{k}{m}\}$ is the graph of the continuous product function $\profunc{g}: I\rightarrow\mathbb{R}^m$, where $\profunc{g}\in\mathscr{T}$. \label{graphofcpf}
	\end{thm}

	\begin{proof}
		Let us define the RB-operator $T:\mathscr{T}\rightarrow\mathscr{T}$ as follows: 
		\begin{equation}
			T(\profunc{f})(\vec{x})=\prodifs{F}\left(\prodifs{L^{-1}}(\vec{x}),\profunc{f}(\prodifs{L^{-1}}(\vec{x}))\right),\label{Rbop}
		\end{equation}
		where $\vec{x}\in\prodifs{I}$. Now, we prove that $T$ is well-defined. Since from the definition of $\prodifs{L}$, we get $$\prodifs{L^{-1}}(\vec{x})=\left(L_{1,i_1}^{-1}(x_1),L_{2,i_2}^{-1}(x_2),\ldots,L_{m,i_m}^{-1}(x_m)\right)\quad\mbox{for all}~ \myvec{\vec{x}}{x}{m}\in\prodifs{I}.$$
		Therefore,
		\begin{align*}
			T(\profunc{f})(\vec{x})=&\prodifs{F}\left(\prod_{k=1}^{m}L_{k,i_k}^{-1}(x_k),\prod_{k=1}^{m}f_k(L_{k,i_k}^{-1}(x_k))\right)\\
			=&\prod_{k=1}^{m}F_{k,i_k}\left(L_{k,i_k}^{-1}(x_k),f_k(L_{k,i_k}^{-1}(x_k))\right),\quad[\mbox{using (\ref{prod of F}})].
		\end{align*}
		Let $\myvec{\vec{x}}{x}{m}\in\partial\prodifs{I}$, where $\partial\prodifs{I}$, denotes the boundary of $\prodifs{I}$. Now, $\vec{x}$ may be anywhere on the boundary $\partial\prodifs{I}$ but we  deal with the case, where it is an intersection point. Let $\vec{x}\in I_{i_1i_2\cdots i_p\cdots,i_m}\cap I_{i_1i_2\cdots i_p+1\cdots,i_m}\cap I_{i_1i_2\cdots i_q\cdots,i_m}\cap I_{i_1i_2\cdots i_q+1\cdots,i_m}\cap\cdots$(up to a finite number of intersection), where $i_p\in\{1,2,\ldots,N_p-1\}$ and $i_q\in\{1,2,\ldots,N_q-1\}$. For simplicity, let $x_p\in I_{p,i_p}\cap I_{p,i_p+1}$ and $x_q\in I_{q,i_q}\cap I_{q,i_q+1}$. This is possible if only if $x_p=x_{p,i_p}$ and $x_q=x_{q,i_q}$. Without loss of generality, we assume that $p<q$. Here we  deal with the following four possible cases.\\ 
		
		Case 1. If $x_{p,i_p}\in I_{p.i_p}$ and $x_{q,i_q}\in I_{q,i_q}$, then
		\begin{align*}
			T(\profunc{f}&)(\vec{x})=\left(F_{1,i_1}(L_{1,i_1}^{-1}(x_1),f_1(L_{1,i_1}^{-1}(x_1))),\ldots,F_{p,i_p}(L_{p,i_p}^{-1}(x_{p,i_p}),f_p(L_{p,i_p}^{-1}(x_{p,i_p}))),\right.\\ &\left.\ldots,F_{q,i_q}(L_{q,i_q}^{-1}(x_{q,i_q}),f_q(L_{q,i_q}^{-1}(x_{q,i_q}))),\ldots,F_{m,i_m}(L_{m,i_m}^{-1}(x_m),f_m(L_{m,i_m}^{-1}(x_m)))\right)\\
			&=\left(F_{1,i_1}\left(L_{1,i_1}^{-1}(x_1),f_1(L_{1,i_1}^{-1}(x_1))\right),\ldots,y_{p,i_p},\right.\\
			&\left.\ldots,y_{q,i_q},\ldots,F_{m,i_m}\left(L_{m,i_m}^{-1}(x_m),f_m(L_{m,i_m}^{-1}(x_m))\right)\right).
		\end{align*}  
		Case 2. If $x_{p,i_p}\in I_{p.i_p}$ and $x_{q,i_q}\in I_{q,i_q+1}$, then   
		\begin{align*}
			T(\profunc{f}&)(\vec{x})=\left(F_{1,i_1}(L_{1,i_1}^{-1}(x_1),f_1(L_{1,i_1}^{-1}(x_1))),\ldots,F_{p,i_p}(L_{p,i_p}^{-1}(x_{p,i_p}),f_p(L_{p,i_p}^{-1}(x_{p,i_p}))),\right.\\ &\left.\ldots,F_{q,i_q+1}(L_{q,i_q+1}^{-1}(x_{q,i_q}),f_q(L_{q,i_q+1}^{-1}(x_{q,i_q}))),\ldots,F_{m,i_m}(L_{m,i_m}^{-1}(x_m),f_m(L_{m,i_m}^{-1}(x_m)))\right)\\
			&=\left(F_{1,i_1}\left(L_{1,i_1}^{-1}(x_1),f_1(L_{1,i_1}^{-1}(x_1))\right),\ldots,y_{p,i_p},\right.\\
			&\left.\ldots,y_{q,i_q},\ldots,F_{m,i_m}\left(L_{m,i_m}^{-1}(x_m),f_m(L_{m,i_m}^{-1}(x_m))\right)\right).
		\end{align*} 
		Case 3. If $x_{p,i_p}\in I_{p.i_p+1}$ and $x_{q,i_q}\in I_{q,i_q}$, then
		\begin{align*}
			T(\profunc{f}&)(\vec{x})=\left(F_{1,i_1}(L_{1,i_1}^{-1}(x_1),f_1(L_{1,i_1}^{-1}(x_1))),\ldots,F_{p,i_p+1}(L_{p,i_p+1}^{-1}(x_{p,i_p}),f_p(L_{p,i_p+1}^{-1}(x_{p,i_p}))),\right.\\ &\left.\ldots,F_{q,i_q}(L_{q,i_q}^{-1}(x_{q,i_q}),f_q(L_{q,i_q}^{-1}(x_{q,i_q}))),\ldots,F_{m,i_m}(L_{m,i_m}^{-1}(x_m),f_m(L_{m,i_m}^{-1}(x_m)))\right)\\
			&=\left(F_{1,i_1}\left(L_{1,i_1}^{-1}(x_1),f_1(L_{1,i_1}^{-1}(x_1))\right),\ldots,y_{p,i_p},\right.\\
			&\left.\ldots,y_{q,i_q},\ldots,F_{m,i_m}\left(L_{m,i_m}^{-1}(x_m),f_m(L_{m,i_m}^{-1}(x_m))\right)\right).
		\end{align*}   
		Case 4. If $x_{p,i_p}\in I_{p.i_p+1}$ and $x_{q,i_q}\in I_{q,i_q+1}$, then 
		\begin{align*}
			T(\profunc{f}&)(\vec{x})=\left(F_{1,i_1}(L_{1,i_1}^{-1}(x_1),f_1(L_{1,i_1}^{-1}(x_1))),\ldots,F_{p,i_p+1}(L_{p,i_p+1}^{-1}(x_{p,i_p}),f_p(L_{p,i_p+1}^{-1}(x_{p,i_p}))),\right.\\ &\left.\ldots,F_{q,i_q+1}(L_{q,i_q+1}^{-1}(x_{q,i_q}),f_q(L_{q,i_q+1}^{-1}(x_{q,i_q}))),\ldots,F_{m,i_m}(L_{m,i_m}^{-1}(x_m),f_m(L_{m,i_m}^{-1}(x_m)))\right)\\
			&=\left(F_{1,i_1}\left(L_{1,i_1}^{-1}(x_1),f_1(L_{1,i_1}^{-1}(x_1))\right),\ldots,y_{p,i_p},\right.\\
			&\left.\ldots,y_{q,i_q},\ldots,F_{m,i_m}\left(L_{m,i_m}^{-1}(x_m),f_m(L_{m,i_m}^{-1}(x_m))\right)\right).
		\end{align*}
		From above, we see that the values of $T(\profunc{f})(\vec{x})$ are same in all the four cases. Similarly, all other possibilities can be checked to conclude that $T(\profunc{f})(\vec{x})$ is well
		defined on the boundary of $\prodifs{I}$. Also, from (\ref{prodrbop}), we have
		\begin{align*}
			T(\profunc{f})(\vec{x})=\prod_{k=1}^{m}F_{k,i_k}\left(L_{k,i_k}^{-1}(x_k),f_k(L_{k,i_k}^{-1}(x_k))\right)=\prod_{k=1}^{m}(T_{k}f_k)(x_k).
		\end{align*} 
		Therefore,
		\begin{equation}
			T(\profunc{f})=T_1f_1\times T_2f_2\times\cdots\times T_mf_m.
			\label{pofrboper}
		\end{equation}
		Hence $T$ is well-defined. Also, if $\vec{x}=(\prodpoint{x}{i})\in \prodifs{I}$, then
		\begin{align*}
			T(\profunc{f})(\vec{x})=&\prod_{k=1}^{m}F_{k,i_k}\left(L_{k,i_k}^{-1}(x_{k,i_k}),f_k(L_{k,i_k}^{-1}(x_{k,i_k}))\right)=\prod_{k=1}^{m}F_{k,i_k}\left(x_{k,0},f_k(x_{k,0})\right)\\
			=&\prod_{k=1}^{m}F_{k,i_k}\left(x_{k,0},y_{k,0}\right)=\prod_{k=1}^{m}y_{k,i_k}.
		\end{align*}
		So, $T(\profunc{f})$ interpolates data points  $$\{((\prodpoint{x}{i}),(\prodpoint{y}{i}))\in I\times\mathbb{R}^m:~\myseries{1}{i_k}{N_k};\myseries{1}{k}{m}\}.$$
		Next, we prove that $T$ is a contraction map on $\mathscr{T}$.
		For that, let $\profunc{f},\profunc{f'}\in\mathscr{T}$, and $\myvec{\vec{x}}{x}{m}\in\prodifs{I}$. Then from (\ref{cont.of F}) and (\ref{Rbop}), we get
		\begin{align*}
			\norm{T(\profunc{f})&(\vec{x})-T(\profunc{f'})(\vec{x})}{2}\\ &\leq|\prodifs{b}|_{\infty}\norm{\profunc{f}(\prodifs{L^{-1}}(\vec{x}))-\profunc{f'}(\prodifs{L^{-1}}(\vec{x}))}{2}\\
			&\leq |\prodifs{b}|_{\infty}\norm{\profunc{f}-\profunc{f'}}{\infty}\\
			&\leq b \norm{\profunc{f}-\profunc{f'}}{\infty},
		\end{align*}
		where  $b=\max\{|\prodifs{b}|_{\infty}:~\myseries{1}{i_k}{N_k};\myseries{1}{k}{m}\}<1$. Now, by taking supremum over all $\vec{x}\in I$, we get
		$$\norm{T(\profunc{f})-T(\profunc{f'})}{\infty}\leq b \norm{\profunc{f}-\profunc{f'}}{\infty}.$$
		This shows that $T$ is a contraction map on the complete space $\mathscr{T}$. Hence by Banach fixed point theorem, $T$ has an unique fixed point $\profunc{g}$ in $\mathscr{T}$. Finally, we see that $G$ is the graph of the function $\profunc{g}$. Let $\tilde{G}=\left\{\left(\vec{x},\profunc{g}(\vec{x})\right):~\vec{x}\in I\right\}$. Then
		\begin{align*}
			&\bigcup\left\{\prodifs{W}(\tilde{G}):~\myseries{1}{i_k}{N_k};\myseries{1}{k}{m}\right\}\\
			&=\bigcup \left\{\prodifs{W}((\vec{x},\profunc{g}(\vec{x})):~\vec{x}\in I,~\myseries{1}{i_k}{N_k};\myseries{1}{k}{m} \right\}\\
			&=\bigcup\left\{\left(\prodifs{L}(\vec{x}),\prodifs{F}(\vec{x},\profunc{g}(\vec{x}))\right):~\vec{x}\in I,~\myseries{1}{i_k}{N_k};\myseries{1}{k}{m}\right\}.
		\end{align*}
		Now, since $\profunc{g}$ is the fixed point of $T$. Hence from (\ref{Rbop}), we get 
		$$\profunc{g}(\prodifs{L}(\vec{x}))=\prodifs{F}(\vec{x},\profunc{g}(\vec{x})).$$
		Therefore,
		\begin{align*}
			&\bigcup\left\{\prodifs{W}(\tilde{G}):~\myseries{1}{i_k}{N_k};\myseries{1}{k}{m}\right\}\\
			&=\bigcup\left\{\left(\prodifs{L}(\vec{x}),\profunc{g}(\prodifs{L}(\vec{x}))\right):~\vec{x}\in I,~\myseries{1}{i_k}{N_k};\myseries{1}{k}{m}\right\}\\	
			&=\left\{\left(\vec{x},\profunc{g}(\vec{x})\right):~\vec{x}\in I\right\}\\
			&=\tilde{G}.
		\end{align*}
		Hence by uniqueness of the attractor, we get $G=\tilde{G}$.
	\end{proof}

\end{document}
